\section{Two clusters, depth-first search, and entropy}
\label{sec:dfs-entropy}

The interpolation argument will join clusters along a short path while
controlling the change in probability. We prepare its three inputs: a
bound on adjacent large distinct clusters, a second-moment estimate for
depth-first pivotal edges, and an entropy bound for several explorations
of the same configuration.

\subsection{Adjacent large distinct clusters}

For an edge $e=xy$ and an integer $s\ge1$, write
\[
 H_s(e)=\{C_x\ne C_y,\ |C_x|\ge s,\ |C_y|\ge s\}.
\]
Cluster size means the number of vertices. The next estimate is uniform
over the edges, including those in different automorphism orbits.
Its proof adapts the cluster-boundary discrepancy argument of
Aizenman, Kesten, and Newman \cite[Section~3]{AizenmanKestenNewman1987}
and Hutchcroft's ghost-field and martingale formulation
\cite[Section~3]{Hutchcroft2020Locality}.
We give the finite-graph averaging argument and its quasi-transitive
passage explicitly.

\begin{proposition}\label{prop:two-cluster}
Let $G$ be an infinite connected locally finite quasi-transitive graph
of subexponential growth. For every $p_*>0$ there is $C<\infty$ such
that
\[
 \PP_p(H_s(e))\le C s^{-1/2}
 \qquad(e\in E(G),\ s\ge1,\ p_*\le p<1).
\]
The constant may depend on $G$ and $p_*$.
\end{proposition}

\begin{proof}
Let $D$ bound the degrees of $G$. We first prove an averaged estimate
on any finite induced subgraph $F$. Independently of the edge states,
mark each vertex with probability $1-e^{-h}$, where $0<h\le1$.
Let $T_e$ be the event that the endpoints of $e$ lie in distinct clusters
both containing a mark. Let $N_e$ count the marked clusters incident
to $e$, counting a cluster only once even when it contains both endpoints,
and put $I_e=\mathbf1_{\{N_e\ge1\}}$. Flipping $e$ does not change the
union of its endpoint clusters, so $I_e$ is unchanged by that flip.
For closed $e$, $\mathbf1_{T_e}=N_e-I_e$, whereas for open $e$,
$N_e=I_e$. With $q=1-p$, conditioning on all other edges therefore gives
\begin{equation}\label{eq:marked-cluster-discrepancy}
 \sum_{e\in E(F)}\PP_p(T_e)
 =\EE_p\sum_{K\text{ marked cluster}}
 \left(c(K)-\frac qp\,o(K)\right),
\end{equation}
where $c(K)$ and $o(K)$ count the closed and open edges incident to $K$,
respectively. Each incidence is counted once per cluster.

Explore a cluster $C_v$, querying all its incident edges, including
internal edges, once each. If $T$ is the number of queries and
$m=|C_v|$, then $T\le Dm$ and, when $T>0$, $m\le T+1\le2T$.
The fresh query bits are independent Bernoulli variables of parameter
$p$; after stopping, extend them with independent auxiliary bits.
Let $Z_j$ be their centered partial sums. Since
$o(C_v)$ is the number of open query bits,
\[
 c(C_v)-\frac qp o(C_v)=-\frac1p Z_T.
\]
Rewrite the sum over clusters in
\eqref{eq:marked-cluster-discrepancy} as a sum over vertices, dividing
the summand by the cluster size. Averaging the marks and taking absolute
values gives
\begin{equation}\label{eq:marked-average-bound}
 \frac1{|V(F)|}\sum_e\PP_p(T_e)
 \le\frac Dp\max_{v\in V(F)}
 \EE_p\left[\frac{|Z_T|}{T}\min\{2hT,1\}\right],
\end{equation}
with value zero when $T=0$.
Doob's maximal inequality yields
\[
 \EE_p\max_{j\le 2^{k+1}}|Z_j|
 \le2\sqrt{pq\,2^{k+1}}.
\]
Partitioning according to $2^k\le T<2^{k+1}$, the expectation in
\eqref{eq:marked-average-bound} is at most
\[
 C_0\sqrt{pq}\sum_{k\ge0}2^{-k/2}
                       \min\{2^{k+2}h,1\}
 \le C_1\sqrt{pq}\sqrt h.
\]
For the last bound, split the sum where $2^kh$ first exceeds one;
both resulting geometric series are bounded by a constant times
$\sqrt h$. On $H_s(e)$, conditional on the edges, the two clusters
are marked with probability at least $(1-e^{-hs})^2$. Taking $h=1/s$
therefore proves
\begin{equation}\label{eq:finite-two-cluster-average}
 \sum_{e\in E(F)}\PP_p^F(H_s(e))
 \le C_2|V(F)|s^{-1/2},
\end{equation}
uniformly for $p\ge p_*$.

To pass to an individual edge, choose $C_0'$ so that every vertex orbit
is $C_0'$-dense. Fix an edge $e=xy$ and a center $o$. For every vertex
of $B(o,n)$, choose a point in the orbit of $x$ within distance $C_0'$.
There are at least $|B(o,n)|/V_+(C_0')$ distinct chosen points. At each,
an automorphic copy of $e$ has both endpoints in $B(o,n+C_0'+1)$.
Each unoriented edge is counted at most twice, so the number of distinct
copies is at least
\[
 \frac{|B(o,n)|}{2V_+(C_0')}.
\]
For fixed $s$, subexponential growth gives arbitrarily large $n$ with
\[
 |B(o,n+C_0'+s+1)|\le2|B(o,n)|.
\]
Otherwise iteration of this fixed-width expansion would force
exponential growth. Set $F=G[B(o,n+C_0'+s+1)]$. A connected cluster
with at least $s$ vertices has a connected $s$-vertex witness within
distance $s-1$ of any specified vertex. Thus $H_s$ in the full graph
implies $H_s$ in $F$ for each selected copy; restriction also preserves
distinctness. All copies have the same full-graph probability.
Apply \eqref{eq:finite-two-cluster-average} and the displayed bounds
to obtain the proposition. The constants do not depend on the orbit
of $e$.
\end{proof}

\subsection{Clusters conditioned to avoid one another}

Conditional association of connection events was established by
van den Berg and Kahn \cite{VanDenBergKahn2001}.
The pivotal estimate uses the cluster-function correlation inequality of
van den Berg, H\"aggstr\"om, and Kahn
\cite[Theorems~1.3 and~1.4]{VanDenBergHaggstromKahn2006}. We include the
product-measure proof by alternating cluster updates, the method of
\cite[Section~2.1]{VanDenBergHaggstromKahn2006}.
Here a cluster object records its vertices and all its open edges.
For a root set, take the union of the cluster objects of its roots.
Objects are ordered by inclusion of both vertices and edges.
A random object is \emph{associated} if any two increasing real-valued
functions of it have nonnegative covariance.

\begin{lemma}[Avoidance correlation]\label{lem:avoidance}
Let $G$ be a finite undirected graph with independent edge probabilities
strictly between zero and one. Let $A,B$ be disjoint nonempty vertex
sets. Conditional on $A\not\leftrightarrow B$, each of the two cluster
objects $C(A)$ and $C(B)$ is associated. If $f$ and $g$ are increasing
functions of the respective objects, then
\[
 \Cov\bigl(f(C(A)),g(C(B))\mid A\not\leftrightarrow B\bigr)\le0.
\]
\end{lemma}

\begin{proof}
Given one exact object under avoidance, the other has its ordinary
product-cluster law after deleting the first object's vertices.
Indeed, the exact object fixes its open internal edges, its closed
internal edges, and its closed boundary edges; all edges wholly outside
it retain their independent laws. Avoidance is then automatic.
By restriction of a common configuration, this conditional law is
stochastically decreasing as the deleted object grows.

Start with a deterministic admissible object at $A$, sample the object
at $B$ by this conditional law, then update the object at $A$ in the
same manner, and repeat. Association is preserved at every half-step.
To see this, conditional on the current object, the next object is an
increasing function of independent bits and is associated by Harris's
inequality. The two conditional means of increasing functions of the
next object are decreasing functions of the current object. Their
covariance is nonnegative when the current object is associated.
The conditional covariance identity proves the assertion.

The two-step chain on the $A$ objects has the avoidance marginal as
a stationary law, because its updates are the two conditional laws of
that joint distribution. Every transition has a common positive mass
at the minimal object consisting of $A$ and no open edges: close all
edges incident to $A$ in its update. This probability is bounded below
by the product of the closed-edge probabilities over all edges of $G$.
Subtracting that common mass from every transition shows that total
variation distances contract by a factor strictly less than one.
Consequently the chain converges to its stationary law. Association
passes to this finite-state limit, proving association of $C(A)$;
the same argument applies to $C(B)$.

Finally, the conditional mean of $g(C(B))$ given $C(A)$ is decreasing
in $C(A)$. Association of the latter object makes its covariance with
the increasing function $f(C(A))$ nonpositive. Taking conditional
expectations proves the claimed cross covariance.
\end{proof}

\subsection{A depth-first pivotal second moment}

Fix a finite undirected graph, a vertex $w$, and a deterministic target
set $B$. Write $E=\{w\leftrightarrow B\}$ and let $\mu$ be homogeneous
Bernoulli percolation with parameter $0<p<1$, with $q=1-p$.
Parallel edges are allowed. We use the following deterministic lazy
depth-first search. On reaching a vertex, push its unqueried edges to
unreached vertices onto a stack in a fixed order. Pop the last entry;
if its head has meanwhile been reached, skip it without reading its
bit. Otherwise query the bit, reaching its head when it is open.
Stop on reaching $B$ or exhausting the stack.

A \emph{query prefix} records the previously queried edges and their
states. At a query of $e$, put
\[
 h=\mu(E\mid\text{prefix}),\qquad
 h_b=\mu(E\mid\text{prefix},X_e=b),\qquad d_e=h_1-h_0.
\]
These probabilities condition only on that search's own query prefix.
They are not conditioned on any additional cluster information that
another construction may possess. Since $E$ is increasing, $0\le d_e\le1$.
A queried edge is an \emph{actual open pivot} if it is open and closing
it destroys $E$ in the full configuration, including unqueried edges.

\begin{lemma}\label{lem:dfs}
For the preceding search, let
\[
 W=\sum_{e\text{ queried actual open pivot}}q d_e.
\]
Then $\EE_\mu W\le1$ and $\EE_\mu W^2\le3$.
\end{lemma}

\begin{proof}
The claim is immediate if $w\in B$ or $B=\varnothing$.
The conditional probabilities of $E$ form a martingale along the
search. At a query its conditional variance increment is $pq d_e^2$.
The search decides $E$, so
\begin{equation}\label{eq:dfs-variance-budget}
 \EE_\mu\sum_{e\text{ queried}}pq d_e^2
 =\Var_\mu(\mathbf1_E)\le1.
\end{equation}
Given the pre-query prefix, the probability of an actual open pivot is
$p d_e$. Thus \eqref{eq:dfs-variance-budget} equals $\EE W$.
It also bounds the diagonal terms in $\EE W^2$, since
$p d_e(qd_e)^2\le pq d_e^2$.

It remains to bound the products of weights from two pivots. We prove
the following conditional estimate for each earlier query $f$:
\begin{equation}\label{eq:dfs-future-pair-bound}
 \EE\left[
 \mathbf1_{\{f\text{ actual pivot}\}}
 \sum_{e\text{ later actual open pivot}}q d_e
 \,\middle|\,\text{prefix of }f,\ f\text{ open}\right]\le d_f.
\end{equation}
Its right side is the pivotal probability before reading $f$.
Multiplying by the probability $p$ that $f$ is open and by its weight
$qd_f$ will bring the pair sum back within the variance budget
\eqref{eq:dfs-variance-budget}.

Fix this prefix and condition on $f$ being open, with new head $z$.
If $z\in B$ there are no later queries. Otherwise contract the vertices
reached before $f$ to a root $A$ and delete $f$ for the connectivity
calculations that follow. Those vertices are connected by known open
edges and contain no target vertex. If $f$ is an actual pivot, $z$
connects to $B$ without using $A$, while $A$ does not connect to $B$.
The depth-first subsearch from $z$ then succeeds before any older stack
entry is resumed. We call this subsearch, stopped on success or
exhaustion, the \emph{phase}.

At a query $e$ in the phase, contract its already reached vertices to
$S$, containing $z$, and wire $B$ to one target vertex. Keep $A$ and
$S$ separate: their only known open link before deletion was $f$.
Delete queried closed edges. All unqueried edges remain independent,
including older stack edges whose endpoints may now lie in $A$ and
$S$. Whether still pending or already skipped, they have not been tested.
The tail of $e$ lies
in $S$ and its head $b$ is unreached.

In this residual product graph with $e$ held closed, let
\[
 \begin{array}{c|cccc}
 \text{connectivity of }A,S,B
 &\text{all separate}&SB\text{ only}&AB\text{ only}&AS\text{ only}\\ \hline
 \text{probability}&u&v&v_A&x.
 \end{array}
\]
The remaining possibility is that all three are connected. Define
\[
 \begin{aligned}
 \alpha&=\PP(\text{all separate},\ b\leftrightarrow B),&
 \beta&=\PP(\text{all separate},\ b\leftrightarrow A),\\
 \gamma&=\PP(SB\text{ only},\ b\leftrightarrow A),&
 \sigma&=\PP(AS\text{ only},\ b\leftrightarrow B).
 \end{aligned}
\]
Here and until the next conditioning all probabilities refer to the
unqueried product bits with $f,e$ deleted. In the actual search $f$
is open, so $A$ and $S$ are wired together when computing $d_e$.
Consequently
\begin{equation}\label{eq:dfs-partition-pivots}
 d_e=\alpha+\sigma,
 \qquad
 \PP(f,e\text{ both actual open pivots}\mid\text{prefix},f\text{ open})
 =p\alpha.
\end{equation}
Indeed, opening $e$ joins the wired source to the target exactly in
the $\alpha$ or $\sigma$ cases. In the second case $A$ was already
connected to $S$ without $f$, so $f$ is not pivotal; the first case
gives both pivots.

Apply Lemma~\ref{lem:avoidance} with root sets $A$ and $S\cup B$,
using the increasing events $b\leftrightarrow A$ and
$S\leftrightarrow B$ in their respective cluster objects. It gives
$\gamma u\le v\beta$. Reversing the roles, with roots $B$ and
$A\cup S$, gives $\sigma u\le x\alpha$. The two decreasing isolation
events $A\not\leftrightarrow S\cup B$ and
$B\not\leftrightarrow A\cup S$ have intersection of probability $u$.
Harris's inequality therefore gives
\begin{equation}\label{eq:dfs-isolation-inequalities}
 \gamma u\le v\beta,
 \qquad \sigma u\le x\alpha,
 \qquad u\ge(u+v)(u+x).
\end{equation}
In particular, when $u>0$,
\begin{equation}\label{eq:dfs-pivotal-comparison}
 d_e\le\alpha\frac{u+x}{u}\le\frac\alpha{u+v}.
\end{equation}
The open-edge information in the cluster objects matters here:
the union of the root vertices alone does not determine
$S\leftrightarrow B$.

We now pay for these later pivotal weights with a second martingale.
An actual pivot $f$ requires the old root $A$ to avoid both $z$ and the
target. Accordingly, for the rest of the phase condition the deleted-$f$
configuration on the fixed event
\[
 U=\{A\not\leftrightarrow\{z,B\}\}.
\]
Let $U_{\rm init}$ be its probability at the beginning of the phase.
If it is zero, no pair involving an actual pivot $f$ contributes.
Otherwise, under this conditioned law, let $r$ be the conditional
probability of $z\leftrightarrow B$ along the phase. This is a
bounded martingale even though its query bits are no longer Bernoulli
with parameter $p$.

At the current prefix write $U_0,U_1$ for the probabilities of $U$
when $e$ is closed and open. The partition above gives
\[
 U_0=u+v,\qquad U_1=U_0-\beta-\gamma,
 \qquad U_{\rm pre}=qU_0+pU_1.
\]
The corresponding probabilities of $U\cap\{z\leftrightarrow B\}$
are $v$ and $v+\alpha-\gamma$. Hence, whenever $U_1>0$,
\begin{equation}\label{eq:dfs-tilted-increment}
 r_1-r_0
 =\frac{U_0\alpha+v\beta-u\gamma}{U_0U_1}
 \ge\frac\alpha{U_1}.
\end{equation}
The conditional bit probabilities under the tilt are
$pU_1/U_{\rm pre}$ and $qU_0/U_{\rm pre}$.

Bayes's formula multiplies the ordinary probability of this prefix by
$U_{\rm pre}/U_{\rm init}$. Thus, after multiplying the tilted
martingale variance budget by $U_{\rm init}$, its contribution at
this prefix, expressed under the ordinary phase law, is
\[
 pq\frac{U_0U_1}{U_{\rm pre}}(r_1-r_0)^2.
\]
If $\alpha>0$, then $u>0$ and $U_1\ge\alpha>0$; moreover
$U_1,U_{\rm pre}\le U_0$. Equations
\eqref{eq:dfs-tilted-increment} and
\eqref{eq:dfs-pivotal-comparison} therefore give
\begin{equation}\label{eq:dfs-converted-budget}
 pq\frac{U_0U_1}{U_{\rm pre}}(r_1-r_0)^2
 \ge pq\frac{U_0\alpha^2}{U_{\rm pre}U_1}
 \ge pq\frac{\alpha^2}{U_0}
 \ge pq\alpha d_e.
\end{equation}
When $\alpha=0$, discard the nonnegative variance contribution.
The total converted variance is at most
\[
 U_{\rm init}r_{\rm init}
 =\PP(U,\ z\leftrightarrow B\mid\text{initial prefix})=d_f.
\]
The last equality is precisely the pivotal event for $f$ before its
bit is read. Using \eqref{eq:dfs-partition-pivots}, the converted variance therefore
bounds the conditional expected sum on the left of
\eqref{eq:dfs-future-pair-bound}. All contributions when $f$ is pivotal
occur inside the phase, as noted above, so this proves that estimate
with the sum over every later query.

Multiply \eqref{eq:dfs-future-pair-bound} by $pq d_f$, average over
the prefix of $f$, and sum over $f$. The sum of products over ordered
pivot pairs is bounded by
$\EE\sum_f pq d_f^2\le1$. Combining this with the diagonal bound
and the factor two for cross terms in $W^2$ proves $\EE W^2\le3$.
\end{proof}

\subsection{Entropy of overlapping transcripts}

For fixed coordinate subsets, related bounds are given by Shearer's
inequality \cite{ChungGrahamFranklShearer1986} and its relative-entropy
forms \cite[Section~IX]{MadimanTetali2010}. The read sets below are
adaptive, and their overlap is bounded only under a possibly dependent
comparison law.
We prove the required estimate by ordering the algorithms randomly
and charging each freshly queried bit to a common transcript.

A transcript records an algorithm's ordered queries and answers,
including its stopping decision. The next result charges several such
transcripts to the information in one common configuration. Its
overlap hypothesis is required only under the law being estimated.

For probability laws $Q,R$ on a finite set, let
\[
 D(Q\Vert R)=\sum_x Q(x)\log\frac{Q(x)}{R(x)}
\]
be their relative entropy, with zero summands omitted.

\begin{lemma}\label{lem:entropy}
Let $\mu$ be the Bernoulli product law of parameter $0<p<1$ on
a finite set of bits, and let $Q$ be any probability law on those bits.
Consider finitely many terminating deterministic adaptive query algorithms, each
reading any bit at most once. Let $Q_i,\mu_i$ be their transcript
laws. Suppose that, on the support of $Q$, every bit is read by at
most $a$ algorithms, where $a\ge1$ is an integer. Then
\[
 \sum_iD(Q_i\Vert\mu_i)
 \le\frac{a^2}{2p(1-p)}D(Q\Vert\mu).
\]
\end{lemma}

\begin{proof}
Write $q=1-p$. At a pre-query node $h$ of algorithm $i$, let $e$
be the bit it queries and put
$b_h=\EE_Q[X_e-p\mid h]$. The transcript entropy chain rule and
the elementary Bernoulli bound
\[
 d(r\Vert p):=r\log\frac rp+(1-r)\log\frac{1-r}q
 \le\frac{(r-p)^2}{pq}
\]
give
\begin{equation}\label{eq:local-transcript-entropy}
 \sum_iD(Q_i\Vert\mu_i)
 \le\frac1{pq}\sum_{i,h}Q(h)b_h^2.
\end{equation}
The Bernoulli reference at each node is still $p$, since its own
history has not read that bit. The displayed bound follows from
$\log t\le t-1$, including endpoint values by continuity.

Assign independent uniform priorities in $[0,1]$ to the algorithms
and run them in priority order, caching previously read bits. Each
algorithm follows exactly its original path; a cached answer supplies
the same bit without a fresh query. For every fixed priority order,
the combined fresh-query transcript is a function of the configuration.
Data processing and the entropy chain rule give total fresh-query
entropy at most $D(Q\Vert\mu)$. The binary lower bound
$d(r\Vert p)\ge2(r-p)^2$ follows because its second derivative in
$r$ is $1/(r(1-r))\ge4$, and its value and first derivative vanish
at $r=p$. It implies
\begin{equation}\label{eq:global-transcript-budget}
 \EE_{\rm priorities}
 \sum_{g\text{ global fresh-query node}}Q(g)b_g^2
 \le\tfrac12D(Q\Vert\mu),
\end{equation}
where $b_g$ is the conditional bias given the entire global pre-query
history and the fixed priorities.

Fix a local node $h$ of algorithm $i$ with $Q(h)>0$, querying $e$,
and condition its priority to be $t$. For a configuration $\omega$,
let $K_e(\omega)$ be the number of algorithms that read $e$.
On configurations visiting $h$, $1\le K_e\le a$, $Q$-almost surely.
The read paths are unchanged by caching. Conditional on such a
configuration, the query at $h$ is fresh exactly when all other
$K_e-1$ users have priority larger than $t$, an event of probability
$(1-t)^{K_e-1}$.
The number $K_e$ may be correlated with $X_e$ under $Q$, so this
freshness weight cannot be pulled out of a conditional expectation of
the bias. Retain their dependence by defining
\[
 z_h(t)=\EE_Q[(X_e-p)(1-t)^{K_e-1}\mid h],\qquad
 w_h(t)=\EE_Q[(1-t)^{K_e-1}\mid h].
\]
Pool the global pre-query nodes at which this local node is visited
fresh, averaging the other priorities. Their total probability is
$Q(h)w_h(t)$, and the sum of their probability-weighted biases is
$Q(h)z_h(t)$. This pooling is legitimate before the bit is read:
the current local node and membership in the cache are determined by
the global pre-query history. Cauchy--Schwarz thus bounds their
total squared-bias contribution below by
\begin{equation}\label{eq:entropy-pooled-bias}
 Q(h)\frac{z_h(t)^2}{w_h(t)},
\end{equation}
with value zero when $w_h(t)=0$.
Every global fresh query belongs to exactly one local node, so these
lower bounds can be summed over all nodes.

At priority zero no other algorithm precedes $i$, so $z_h(0)=b_h$:
this endpoint recovers the original local bias. The overlap bound makes
$z_h$ a polynomial of degree at most $a-1$, since
$K_e\in\{1,\ldots,a\}$ on the conditional $Q$ support. Also
$0\le w_h\le1$. For any polynomial $z$ of degree
at most $a-1$, the classical shifted-Legendre endpoint bound
\cite[Tables~18.3.1 and~18.6.1]{DLMF} gives
\begin{equation}\label{eq:polynomial-endpoint}
 |z(0)|^2\le a^2\int_0^1|z(t)|^2\,dt.
\end{equation}
Indeed, expand in the orthonormal shifted Legendre polynomials of
degrees $0,\ldots,a-1$. Their squared endpoint values are
$1,3,\ldots,2a-1$, with sum $a^2$, and Cauchy--Schwarz gives
\eqref{eq:polynomial-endpoint}.
Integrating \eqref{eq:entropy-pooled-bias} therefore yields at least
$Q(h)b_h^2/a^2$. Combine with
\eqref{eq:global-transcript-budget} and then
\eqref{eq:local-transcript-entropy} to prove the lemma.

No overlap bound under $\mu$ was used. The random variable $K_e$ may
be correlated with $X_e$ under $Q$; this correlation is retained inside
$z_h(t)$ throughout the argument.
\end{proof}
